% !TeX encoding = UTF-8
\SetChapterInfo{subtitle={由完备性获得统一控制与连续性},keywords={一致有界; 测度极限; 开映射; 闭图像; 亚椭圆性},prereq={Baire 定理、第一章函数空间与分布},goals={掌握三大基本定理; 区分闭性与有界性; 理解定性正则性如何产生估计}}
\chapter{Baire--Hausdorff 定理的应用}\label{ch:yosida-02}
\ChapterSource{本章依照 \cite{yosida-zh,yosida-en} 第二章七节的顺序重新组织。实际核对范围：中译本印刷第 71--83 页（PDF 第 86--98 页），英文第六版印刷第 68--81 页（PDF 第 83--96 页）。中译本第 84 页为空白页；英文第 81 页下部已进入第三章。以下采用本稿记号与编号，补充步骤另作说明，并非逐字翻译。}

第一章提供了大量完备函数空间。本章解释完备性如何把逐点信息变成统一控制：一致有界性把各点上的有界性变为邻域上的一致估计；开映射把满射变为可控地求解方程；闭图像把极限相容性变为连续性。三者都以第零章的 Baire 定理为基础。

除特别说明外，标量域为 $\R$ 或 $\mathbb C$，所有空间均为 Hausdorff 拓扑向量空间。沿用第一章术语：$F$ 空间的拓扑由完备平移不变度量给出，不要求局部凸；Fr\'echet 空间另要求局部凸。Banach 空间是完备赋范空间。第一纲集是可数个无处稠密集之并，其补集称为余第一纲集。非空完备度量空间中的余第一纲集稠密。

\section{一致有界性定理及共鸣定理}\label{sec:yosida-02-01}
\noindent\textit{原书位置：中译本第 71--72 页；英文版第 68--70 页。}

\begin{theorem}[一致有界性：原书的次可加形式]\label{thm:y2-uniform-general}
设 $X$ 是非第一纲的 Hausdorff 线性拓扑空间，$Y$ 具有第一章意义的拟范数 $q$，$T_a:X\to Y$ 连续，并满足
\[
 q(T_a(x+y))\le q(T_ax)+q(T_ay),\qquad
 q(T_a(tx))=q(tT_ax)\quad(t\ge0).
\]
若对每个 $x\in X$，集合 $\{T_ax:a\in A\}$ 拓扑有界，则对任意 $\varepsilon>0$，存在零邻域 $V$，使 $q(T_ax)\le\varepsilon$ 对所有 $x\in V,a\in A$ 成立。
\end{theorem}
\begin{proof}
固定 $\varepsilon>0$，令
\[
 E_m=\left\{x:\sup_a\bigl[q(m^{-1}T_ax)+q(m^{-1}T_a(-x))\bigr]\le\varepsilon/2\right\}.
\]
各 $E_m$ 闭。拓扑有界性意味着对固定 $x$，$m^{-1}T_a(\pm x)$ 对 $a$ 一致趋零，故 $X=\bigcup_mE_m$。于是某 $E_m$ 含有 $x_0+W$，其中 $W$ 是零邻域且 $0\in W$。对 $h\in W$，次可加性和正数缩放条件给出
\[
 q(T_a(h/m))\le q(m^{-1}T_a(x_0+h))+q(m^{-1}T_a(-x_0))\le\varepsilon.
\]
取 $V=m^{-1}W$ 即可。拟范数不要求绝对齐次，不能把拓扑有界误换成拟范数数值有界。
\end{proof}

\begin{theorem}[Banach--Steinhaus 与共鸣定理]\label{thm:y2-banach-steinhaus}
设 $X$ 为 Banach 空间，$Y$ 为赋范空间，$T_a:X\to Y$ 为有界线性算子。若 $\sup_a\|T_ax\|<\infty$ 对每个 $x\in X$ 成立，则 $\sup_a\|T_a\|<\infty$。更一般地，若 $X$ 是 Baire 局部凸空间、$Y$ 是局部凸空间，则逐点有界的连续线性算子族等度连续。
\end{theorem}
\begin{proof}
令 $E_m=\{x:\sup_a\|T_ax\|\le m\}$。这些闭集覆盖 $X$，故某 $E_m$ 包含 $B(x_0,r)$。当 $\|h\|<r$ 时，$x_0+h,x_0\in E_m$，故 $\sup_a\|T_ah\|\le2m$。对单位向量取 $h=(r/2)x$，得 $\sup_a\|T_a\|\le4m/r$。

一般版本对每个目标连续半范数 $q$ 使用 $E_m=\{x:\sup_aq(T_ax)\le m\}$。同样的平移相减给出某零邻域上的统一界，再缩放得到任意小的 $q$ 界。对有限个目标半范数条件取源邻域的交，即得等度连续性。特别地，Fr\'echet 空间上逐点有界的连续线性泛函族满足 $\sup_a|T_ax|\le C\max_{j\le r}p_j(x)$，其中 $p_j$ 是有限多个源半范数。
\end{proof}

图 \ref{fig:y2-uniform} 对应 Banach 版本中从非零中心的内球转到零附近估计的步骤。
\begin{figure}[htbp]
\centering
\begin{tikzpicture}[fa diagram]
  \draw[fa region] (-2,-1.7) rectangle (2.2,1.8);
  \node at (-1.45,1.4) {$E_m$};
  \draw[fa main] (.2,0) circle (1);
  \node[fa dot,label=below:$x_0$] at (.2,0) {};
  \node[fa dot,label=above:$x_0+h$] at (.8,.45) {};
  \draw[->,fa accent] (.2,0)--(.8,.45);
  \draw[->,fa main] (2.6,.1)--(4.3,.1) node[midway,above] {相减};
  \draw[fa main,fill=LogicColor!8] (6,0) circle (1);
  \node[fa dot,label=below:$0$] at (6,0) {};
  \node[fa dot,label=above:$h$] at (6.6,.45) {};
  \draw[->,fa accent] (6,0)--(6.6,.45);
  \node[fa note] at (2.9,-2.4) {$x_0+h,x_0\in E_m\quad\Longrightarrow\quad\sup_a\|T_ah\|\le2m$};
\end{tikzpicture}
\caption{一致有界性证明的转折：Baire 定理在某个 $E_m$ 中找到非空内球，线性相减把它变为零附近的统一估计，再缩放至单位球。仅取 $x_0$ 与 $x_0+h$ 相减已足够，不要求 $E_m$ 本身是子空间。}
\label{fig:y2-uniform}
\end{figure}



\begin{corollary}[强极限的连续性]\label{cor:y2-strong-limit}
若 $T_n:X\to Y$ 如上且 $T_nx\to Tx$ 对每个 $x$ 成立，则 $T$ 有界线性，且 $\|T\|\le\liminf_n\|T_n\|$。
\end{corollary}
\begin{proof}
收敛序列逐点有界，故 $M=\sup_n\|T_n\|<\infty$。在和与数乘的恒等式中取极限得线性，并有 $\|Tx\|\le M\|x\|$。又由 $\|T_nx\|\le\|T_n\|\|x\|$ 取下极限，再对单位球取上确界即得结论。这里不要求 $Y$ 完备，因为已经假定极限在 $Y$ 中存在。
\end{proof}

\begin{theorem}[Neumann 级数]\label{thm:y2-neumann}
若 $X$ 为 Banach 空间，$A\in\mathcal L(X)$ 且 $\|A\|<1$，则
\[
 (I-A)^{-1}=\sum_{j=0}^\infty A^j,\qquad
 \|(I-A)^{-1}\|\le(1-\|A\|)^{-1}.
\]
级数在算子范数下收敛。原书以 $A=I-T$ 表述此结果。
\end{theorem}
\begin{proof}
第一章已证明 $\mathcal L(X)$ 完备，而 $\sum_j\|A^j\|\le\sum_j\|A\|^j<\infty$。若 $S_N=\sum_{j=0}^NA^j$，则 $(I-A)S_N=S_N(I-A)=I-A^{N+1}$。按算子范数取极限得双侧逆，范数界来自几何级数。
\end{proof}

\begin{example}[源空间完备性的作用]
给有限支集序列空间 $c_{00}$ 配上 $\|x\|_\infty$，令 $T_nx=nx_n$。每个固定 $x$ 只在有限多个坐标非零，所以 $\sup_n|T_nx|<\infty$，但 $\|T_n\|=n$。$c_{00}$ 不完备，其完备化为 $c_0$。
\end{example}

\begin{exercise}
设 $X$ 为 Banach 空间，$T_n\in\mathcal L(X,Y)$，$Y$ 赋范。证明：$T_nx\to0$ 对所有 $x$ 成立，当且仅当 $\sup_n\|T_n\|<\infty$ 且在某个稠密子集 $D\subset X$ 上 $T_nd\to0$。
\end{exercise}
\begin{solution}
必要性由共鸣定理，且可取 $D=X$。反向令 $M=\sup_n\|T_n\|$。给定 $x,\varepsilon$，取 $d\in D$ 使 $M\|x-d\|<\varepsilon/2$（$M=0$ 时显然）。当 $n$ 足够大时 $\|T_nd\|<\varepsilon/2$，于是 $\|T_nx\|\le M\|x-d\|+\|T_nd\|<\varepsilon$。无需 $D$ 为子空间。
\end{solution}

\begin{exercise}
设 $T:X\to Y$ 为 Banach 空间之间的有界可逆算子，$E\in\mathcal L(X,Y)$ 且 $\|T^{-1}E\|<1$。证明 $T+E$ 可逆，给出逆的级数及范数界。
\end{exercise}
\begin{solution}
由 $T+E=T(I+T^{-1}E)$，对 $A=-T^{-1}E$ 使用 Neumann 级数得
\[
 (T+E)^{-1}=\sum_{j=0}^\infty(-T^{-1}E)^jT^{-1},\qquad
 \|(T+E)^{-1}\|\le\frac{\|T^{-1}\|}{1-\|T^{-1}E\|}.
\]
级数在 $\mathcal L(Y,X)$ 中按范数收敛，并给出双侧逆。
\end{solution}

\section{Vitali--Hahn--Saks 定理}\label{sec:yosida-02-02}
\noindent\textit{原书位置：中译本第 73--74 页；英文版第 70--72 页。}

本节测度均为可数可加的有限值实或复测度，$|\nu|$ 表示全变差。对正测度 $m$，这里的 $m$-绝对连续明确指
\[
 \forall\varepsilon>0\ \exists\delta>0:\quad m(E)<\delta\Longrightarrow|\nu(E)|<\varepsilon.
\]
对有限值复测度，它等价于在 $m$-零集上为零。非平凡方向可用反证法选 $m(E_j)<2^{-j}$ 而 $|\nu(E_j)|\ge\varepsilon$，令 $F_k=\bigcup_{j\ge k}E_j$。则 $m(\bigcap_kF_k)=0$，由 $|\nu|$ 的从上连续性得 $|\nu|(F_k)\to0$，矛盾。此论证使用了第一章已证明的复测度全变差有限。

\begin{lemma}[集合空间的完备性]\label{lem:y2-measure-algebra}
在正测度空间 $(S,\Sigma,m)$ 上，令 $\Sigma_0=\{E\in\Sigma:m(E)<\infty\}$，把相差 $m$-零集的集合视为同一点。则 $d([E],[F])=m(E\triangle F)$ 是完备度量。
\end{lemma}
\begin{proof}
由 $\|\mathbf1_E-\mathbf1_F\|_1=m(E\triangle F)$ 得到向 $L^1$ 的等距嵌入。若 $\mathbf1_{E_j}$ 为 Cauchy 序列，$L^1$ 完备性给出极限 $f$，并有几乎处处收敛的子列。子列只取 $0,1$，所以 $f=\mathbf1_E$ 几乎处处；用子列的可测上极限可选择可测 $E$。又 $m(E)=\|f\|_1<\infty$，所以极限仍在像中。
\end{proof}

\begin{theorem}[Vitali--Hahn--Saks]\label{thm:y2-vhs}
设 $m(S)<\infty$，$\nu_n$ 为 $m$-绝对连续的复测度。若对每个 $E\in\Sigma$，有限极限 $\nu(E)=\lim_n\nu_n(E)$ 存在，则
\[
 \forall\varepsilon>0\ \exists\delta>0:\quad
 m(E)<\delta\Longrightarrow\sup_n|\nu_n(E)|<\varepsilon.
\]
而且 $\nu$ 是 $m$-绝对连续的复测度。
\end{theorem}
\begin{proof}
每个 $[E]\mapsto\nu_n(E)$ 在引理的完备空间上良定义且连续，因为
\[
 |\nu_n(E)-\nu_n(F)|\le|\nu_n(E\setminus F)|+|\nu_n(F\setminus E)|
\]
在 $m(E\triangle F)\to0$ 时趋零。固定 $\varepsilon>0$，令
\[
 F_k=\{[E]:|\nu_n(E)-\nu_l(E)|\le\varepsilon/4\ \text{对所有 }n,l\ge k\}.
\]
这些闭集由逐集合 Cauchy 性覆盖整个空间。Baire 定理给出 $k,E_0,\eta>0$，使 $d([E],[E_0])<\eta$ 蕴含 $[E]\in F_k$。若 $m(B)<\eta$，令 $E_1=E_0\cup B$、$E_2=E_0\setminus B$。两者均在该球内，且 $\nu_n(B)=\nu_n(E_1)-\nu_n(E_2)$，故对 $n\ge k$ 有
\[
 |\nu_n(B)-\nu_k(B)|\le\varepsilon/2.
\]
再用有限多个 $\nu_1,\ldots,\nu_k$ 的绝对连续性缩小 $\delta<\eta$，使它们在 $m(B)<\delta$ 时均小于 $\varepsilon/4$。便得全部 $n$ 的统一界小于 $\varepsilon$。

取极限得 $\nu$ 的绝对连续性，逐集合取极限也给出有限可加性。若 $E_j$ 两两不交，令 $E=\bigcup_jE_j$、$R_N=\bigcup_{j>N}E_j$。因 $m(S)<\infty$，$m(R_N)\to0$，故
\[
 \nu(E)-\sum_{j=1}^N\nu(E_j)=\nu(R_N)\longrightarrow0.
\]
这证明可数可加性。有限值复测度的全变差自动有限；此处无需额外假定变差对 $n$ 一致有界，这也对应中译本第 73 页的校注。
\end{proof}

\begin{corollary}[一致可数可加性与无限参考测度]\label{cor:y2-nikodym}
若 $\nu_n(E)$ 对每个可测集收敛，无需预给参考测度，则对 $E_j\downarrow\varnothing$ 有 $\sup_n|\nu_n(E_j)|\to0$，逐集合极限是复测度。若各 $\nu_n$ 对任意正测度 $m$ 满足上述绝对连续性，则极限也对 $m$ 绝对连续，即使 $m(S)=\infty$。
\end{corollary}
\begin{proof}
定义有限正控制测度
\[
 \mu(E)=\sum_{n=1}^\infty2^{-n}\frac{|\nu_n|(E)}{1+|\nu_n|(S)}.
\]
非负求和可交换，故 $\mu$ 可数可加且 $\mu(S)\le1$。每个 $|\nu_n|\le2^n(1+|\nu_n|(S))\mu$，从而各 $\nu_n$ 都对 $\mu$ 绝对连续。应用定理，$\mu(E_j)\to0$ 导出一致从上连续性和极限的可数可加性。分母的 $1+$ 处理了零测度。

若另有 $m$，在 $\{[E]:m(E)<\infty\}$ 上重复定理的闭集和 Baire 论证，便得对 $m$ 的统一绝对连续性；所有 $m(B)<\delta$ 的集合本来就在该空间内。因此极限也对 $m$ 绝对连续。可数可加性已由有限控制测度获得，不需要在无限 $m$ 下错误地断言所有尾集的 $m$ 测度趋零。
\end{proof}

\begin{example}[密度几乎处处收敛不能替代逐集合收敛]
在 $[0,1]$ 上令 $f_n=n\mathbf1_{(0,1/n)}$、$\nu_n(E)=\int_Ef_n\,dx$。虽然 $f_n\to0$ 几乎处处，取 $E_n=(0,1/n)$ 却有 $m(E_n)\to0$、$\nu_n(E_n)=1$。所以该列不一致绝对连续，由定理它不可能对每个可测集都有极限。
\end{example}

\begin{exercise}
证明定理的假设还蕴含：对每个 $\varepsilon>0$ 存在 $\delta>0$，使 $m(E)<\delta$ 时 $\sup_n|\nu_n|(E)<\varepsilon$。
\end{exercise}
\begin{solution}
先证一般估计 $|\rho|(E)\le4\sup_{B\subset E}|\rho(B)|$，其中 $B$ 可测。实带符号测度 $a$ 的正负分解给出 $|a|(E)\le2\sup_{B\subset E}|a(B)|$。写 $\rho=a+ib$，由分割定义 $|\rho|(E)\le|a|(E)+|b|(E)$，即得估计。若 $m(E)<\delta$，所有 $B\subset E$ 也满足此条件。将定理的误差取为 $\varepsilon/8$，得 $\sup_n|\nu_n|(E)\le\varepsilon/2<\varepsilon$。
\end{solution}

\section{广义函数序列的逐项可微性}\label{sec:yosida-02-03}
\noindent\textit{原书位置：中译本第 75 页；英文版第 72 页。}

\begin{theorem}[分布的逐测试函数极限]\label{thm:y2-distribution-limit}
设 $T_n\in\Dist(\Omega)$，对每个 $\varphi\in\Test(\Omega)$，$T(\varphi)=\lim_n\Pair{T_n}{\varphi}$ 存在且有限。则 $T\in\Dist(\Omega)$。对每个 $K\Subset\Omega$，存在 $C_K>0$ 和整数 $m_K\ge0$，使
\[
 \sup_n|\Pair{T_n}{\varphi}|\le C_K
 \max_{|\alpha|\le m_K}\sup_K|\partial^\alpha\varphi|
 \quad(\varphi\in\Test_K).
\]
\end{theorem}
\begin{proof}
第一章已证明 $\Test_K$ 是 Fr\'echet 空间。各 $T_n|_{\Test_K}$ 为连续线性泛函，逐点收敛使它们逐点有界。由定理 \ref{thm:y2-banach-steinhaus} 得等度连续性，再把有限多个导数半范数合并，即得所述估计。取极限得到 $T$ 的同一估计；线性直接由极限得到。对每个紧集的有限阶估计正是第一章的分布连续性判别，所以 $T$ 为分布。这里在各 $\Test_K$ 上使用 Baire 定理，并未假定整个 LF 空间 $\Test(\Omega)$ 是 Fr\'echet 空间。
\end{proof}

\begin{corollary}[逐项求导与光滑乘法]\label{cor:y2-termwise}
若 $T_n\to T$ 是逐测试函数的分布收敛，则对任意多重指标 $\alpha$ 及 $a\in C^\infty(\Omega)$，$\partial^\alpha T_n\to\partial^\alpha T$，且 $aT_n\to aT$。
\end{corollary}
\begin{proof}
对固定 $\varphi$，$\partial^\alpha\varphi,a\varphi\in\Test(\Omega)$，且
\[
 \Pair{\partial^\alpha T_n}{\varphi}
 =(-1)^{|\alpha|}\Pair{T_n}{\partial^\alpha\varphi}
 \longrightarrow(-1)^{|\alpha|}\Pair{T}{\partial^\alpha\varphi}.
\]
乘法同理。该收敛不等于导数作为普通函数逐点或一致收敛。
\end{proof}

\begin{example}[正则化阶跃与 Dirac 质量]
取 $\rho\in\Test(\R)$、$\rho\ge0$、$\int\rho=1$，令 $\rho_\varepsilon(x)=\varepsilon^{-1}\rho(x/\varepsilon)$，$H_\varepsilon(x)=\int_{-\infty}^x\rho_\varepsilon(t)\,dt$。由有界收敛，$H_\varepsilon\to H=\mathbf1_{(0,\infty)}$ 于分布意义；故 $H_\varepsilon'=\rho_\varepsilon\to H'=\delta_0$。求导后极限已不是普通局部可积函数。
\end{example}

\begin{exercise}
令 $u_n(x)=\sin(nx)$。证明 $u_n\to0$、$u_n'\to0$ 于分布意义，而 $u_n'$ 不逐点收敛到零。说明分布级数的逐项求导需要什么假设。
\end{exercise}
\begin{solution}
分部积分得 $|\int\sin(nx)\varphi(x)\,dx|\le n^{-1}\|\varphi'\|_1\to0$。推论给出 $u_n'\to0$，但 $u_n'(0)=n$。若级数部分和 $S_N=\sum_{j=1}^NT_j$ 在每个测试函数上收敛到有限值，则定理保证极限为分布，并且 $\partial^\alpha S_N\to\partial^\alpha S$。无需级数以函数范数绝对收敛；也不能只凭各项可微便断言级数可逐项求导。
\end{solution}

\section{奇点凝聚原理}\label{sec:yosida-02-04}
\noindent\textit{原书位置：中译本第 75--77 页；英文版第 72--75 页。}

\begin{theorem}[有界点集合的二择性]\label{thm:y2-bounded-points}
设 $X$ 为 Banach 空间，$T_j:X\to Y_j$ 为有界线性算子，$Y_j$ 赋范。则
\[
 B=\{x:\sup_j\|T_jx\|<\infty\}
\]
或者等于 $X$，或者是第一纲集。
\end{theorem}
\begin{proof}
写 $B=\bigcup_m E_m$，其中 $E_m=\{x:\sup_j\|T_jx\|\le m\}$ 闭。若 $B$ 不是第一纲集，某 $E_m$ 有内点。与共鸣定理相同的平移相减和缩放给出 $\sup_j\|T_j\|<\infty$，故 $B=X$。反之，各 $E_m$ 均闭且无内点，便是无处稠密集。
\end{proof}

\begin{corollary}[Banach--Steinhaus 奇点凝聚原理]\label{cor:y2-condensation}
设 $T_{p,q}:X\to Y_p$ 为有界线性算子，$p,q\ge1$。若对每个 $p$ 存在 $x_p$ 使 $\sup_q\|T_{p,q}x_p\|=\infty$，则
\[
 G=\{x:\sup_q\|T_{p,q}x\|=\infty\ \text{对所有 }p\}
\]
是稠密的余第一纲集。这里可等价地把上确界无穷写为 $\limsup_{q\to\infty}\|T_{p,q}x\|=\infty$，不能换成普通极限。
\end{corollary}
\begin{proof}
对每个 $p$，定理中的有界点集合 $B_p$ 不等于 $X$，所以是第一纲集。$G=X\setminus\bigcup_pB_p$ 的补集仍是第一纲集，由 Baire 定理得稠密性。
\end{proof}

\begin{theorem}[处处无有限导数的连续函数：补全锯齿论证]\label{thm:y2-nowhere-differentiable}
在 $C([0,1];\R)$ 中，处处无有限导数的函数构成一个包含稠密余第一纲集的集合；端点的导数按单侧理解。
\end{theorem}
\begin{proof}
对 $m\ge1$ 令
\[
 M_m=\left\{f:\exists t\in[0,1/2],\quad
 |f(t+h)-f(t)|\le mh\ \text{对所有 }0<h<1/2\right\}.
\]
若 $f_j\to f$ 一致且 $f_j\in M_m$，由相应 $t_j$ 的紧性可取 $t_j\to t$。固定 $h\in(0,1/2)$，连续性和一致收敛使不等式传到 $t$，故 $f\in M_m$；所以 $M_m$ 闭。

现证 $M_m$ 无内点。给定 $f$ 和误差 $\varepsilon$，先用分段线性 $g$ 一致逼近 $f$ 到 $\varepsilon/2$，记其有限个斜率绝对值的最大值为 $L$。在 $g$ 的每个线性区间上细分，叠加端点为零、振幅小于 $\varepsilon/2$ 的连续三角波 $w$，使每段斜率为 $\pm M$ 且 $M>L+m$；只需把三角波周期取得足够短，并使每个原区间包含整数个周期。于是 $g+w$ 的每个线性小段斜率的绝对值都大于 $m$。对每个 $t\in[0,1/2]$，即使 $t$ 是顶点，也有一个充分短的右侧线性区间，其差商绝对值大于 $m$。所以 $g+w\notin M_m$，而 $\|g+w-f\|_\infty<\varepsilon$。

因此 $G_+=C([0,1])\setminus\bigcup_mM_m$ 是稠密余第一纲集。若 $f$ 在 $t\in[0,1/2]$ 有有限右导数，则小 $h$ 的差商有界；较大 $h$ 的差商由 $2\|f\|_\infty/h$ 控制，故 $f\in M_m$ 对某个 $m$ 成立。由此 $G_+$ 中函数在左半区间无有限右导数。反射 $Rf(t)=f(1-t)$ 是同胚，$G_+\cap R(G_+)$ 仍为稠密余第一纲集，其函数在整个区间无有限导数。
\end{proof}

\begin{example}[连续周期函数的 Fourier 部分和可以无界]\label{ex:y2-fourier-divergence}
在实 Banach 空间 $C(\mathbb T)$ 上使用标准 Dirichlet 核
\[
 D_N(s)=\sum_{k=-N}^Ne^{iks}
 =\frac{\sin((N+1/2)s)}{\sin(s/2)},\qquad
 S_Nf(t)=\frac1{2\pi}\int_{-\pi}^{\pi}f(t-s)D_N(s)\,ds.
\]
对固定 $t$，$f\mapsto S_Nf(t)$ 为有界实线性泛函，其范数为 $(2\pi)^{-1}\|D_N\|_1$。上界来自积分估计；下界可取连续周期函数近似 $\operatorname{sgn}D_N$，例如 $D_N/\sqrt{D_N^2+\delta^2}$，再令 $\delta\downarrow0$，并相应平移反射。

这些范数无界。事实上 $\sin(s/2)\le s/2$ 给出
\[
 \int_0^\pi|D_N(s)|\,ds
 \ge2\int_0^{(N+1/2)\pi}\frac{|\sin v|}{v}\,dv
 \ge c\sum_{j=1}^{N-1}\frac1{j+1},
\]
最后一步在 $[j\pi+\pi/6,j\pi+5\pi/6]$ 上用 $|\sin v|\ge1/2$；$c>0$ 与 $N$ 无关。因此由共鸣定理的逆否命题，对每个固定 $t$ 存在 $f_t$ 使 $\sup_N|S_Nf_t(t)|=\infty$。

给定任意可数点集，加入一个可数稠密集得到 $\{t_p\}$。奇点凝聚原理给出一个连续 $f$，使上述无界性在全部 $t_p$ 同时发生。对这个 $f$，令
\[
 F_m=\bigcap_N\{t:|S_Nf(t)|\le m\}.
\]
各 $F_m$ 闭，且不可能含非空开弧，否则与 $\{t_p\}$ 稠密且部分和无界矛盾。所以 $P=\mathbb T\setminus\bigcup_mF_m$ 是余第一纲集，并包含所有预先指定的点。在 $P$ 上有 $\limsup_N|S_Nf(t)|=\infty$。

补充说明：$P$ 的势为连续统，无需连续统假设。写 $P=\bigcap_mO_m$ 为稠密开集之交，在任意开弧内递归选择两两分离的闭子弧，令第 $m$ 层均落在 $O_1\cap\cdots\cap O_m$，直径趋零。每个二进制无限序列确定一个不同的交点，因而向 $P$ 嵌入 $\{0,1\}^{\mathbb N}$；反向上界由 $P\subset\mathbb T$ 给出。这展开了原书关于无须连续统假设的注记。
\end{example}

\begin{exercise}
设 $X$ 为 Banach 空间，$T_n:X\to Y_n$ 有界线性且 $\sup_n\|T_n\|=\infty$。证明存在稠密余第一纲集 $G$，使每个 $x\in G$ 满足 $\limsup_n\|T_nx\|=\infty$。该结论是否保证 $\|T_nx\|\to\infty$？
\end{exercise}
\begin{solution}
若有界点集合等于 $X$，共鸣定理给出算子范数一致有界，矛盾。因此有界点集合为第一纲集，其补集就是所求 $G$。不保证普通极限：在 $X=Y_n=\R$ 上取 $T_{2k}=0$、$T_{2k-1}x=kx$。对 $x\ne0$ 上极限无穷，而偶数子列恒为零。
\end{solution}

\section{开映射定理}\label{sec:yosida-02-05}
\noindent\textit{原书位置：中译本第 77--79 页；英文版第 75--77 页。}

\begin{lemma}[像的闭包含零邻域]\label{lem:y2-almost-open}
设 $T:X\to Y$ 连续线性，且其值域 $R(T)$ 在 $Y$ 中不是第一纲集。那么对每个 $X$ 的零邻域 $U$，$\overline{T(U)}$ 包含 $Y$ 的某个零邻域。
\end{lemma}
\begin{proof}
取对称零邻域 $W$ 使 $W-W\subset U$。零邻域吸收每个向量，所以 $X=\bigcup_{n\ge1}nW$，$R(T)\subset\bigcup_n n\overline{T(W)}$。若所有这些闭集无内点，值域便是第一纲集，矛盾。故 $\overline{T(W)}$ 有内点。取非空开集 $O\subset\overline{T(W)}$，则 $O-O$ 为零邻域，且
\[
 O-O\subset\overline{T(W)}-\overline{T(W)}
 \subset\overline{T(W-W)}\subset\overline{T(U)}.
\]
中间包含关系由减法连续性得到。
\end{proof}

\begin{theorem}[Banach 开映射定理]\label{thm:banach-open}
设 $X,Y$ 为 $F$ 空间，$T:X\to Y$ 为连续线性满射。则 $T$ 是开映射。特别地，它适用于 Banach 空间和 Fr\'echet 空间。
\end{theorem}
\begin{proof}
取给出拓扑的完备平移不变度量 $d_X,d_Y$，以 $B_X(r),B_Y(r)$ 表示零点开球。$Y$ 完备，故不是第一纲集。由引理，对每个 $r>0$，$\overline{T(B_X(r))}$ 含某零球。

固定 $\varepsilon>0$，令 $r_j=\varepsilon2^{-j-2}$，$j\ge0$。选择 $\eta_j>0$ 使
\[
 B_Y(\eta_j)\subset\overline{T(B_X(r_j))},\qquad
 \eta_j\le2^{-j}.
\]
任取 $y\in B_Y(\eta_0)$。由闭包定义可取 $x_0\in B_X(r_0)$，使残差 $y-Tx_0\in B_Y(\eta_1)$；再取 $x_1\in B_X(r_1)$，使 $y-T(x_0+x_1)\in B_Y(\eta_2)$。如此继续得到
\[
 d_X(x_j,0)<r_j,\qquad
 d_Y\left(y-T\sum_{j=0}^Nx_j,0\right)<\eta_{N+1}.
\]
平移不变度量的三角不等式使部分和为 Cauchy 序列，故在 $X$ 中趋于 $x$，且 $d_X(x,0)\le\sum_jr_j=\varepsilon/2<\varepsilon$。连续性和 $\eta_j\to0$ 给出 $Tx=y$。因此 $B_Y(\eta_0)\subset T(B_X(\varepsilon))$。平移到任意点即得开性。此证明只用度量的三角不等式，未假定度量对标量绝对齐次，也未用局部凸性。
\end{proof}

\begin{corollary}[有界逆与商空间]\label{cor:y2-inverse}
连续线性双射 $T:X\to Y$ 在 $F$ 空间之间具有连续逆。若 $X,Y$ 为 Banach 空间，则逆有界。若 $T$ 只是 Banach 空间之间的有界满射，则存在 $C>0$，对每个 $y$ 可选 $Tx=y$ 且 $\|x\|\le C\|y\|$；诱导映射 $X/\ker T\to Y$ 是 Banach 空间同构。
\end{corollary}
\begin{proof}
逆像开集的条件正是 $T$ 的开性。对满射，由开映射得 $B_Y(\delta)\subset T(B_X(1))$。当 $y\ne0$ 时，将 $y$ 缩放至范数 $\delta/2$ 再选原像，得 $\|x\|<2\|y\|/\delta$；$y=0$ 取 $x=0$。核闭，商空间由第一章知为 Banach，诱导映射为有界双射，再用有界逆结论。这里选取原像的估计并不自动产生一个线性的右逆。
\end{proof}

\begin{example}[等价范数与不完备反例]
同一线性空间上的两个 Banach 范数若满足 $\|x\|_1\le C\|x\|_2$，则反向估计也成立：对恒等双射应用有界逆定理。完备性缺失时可失败。在 $C^1[0,1]$ 上，$\|f\|_{C^1}=\|f\|_\infty+\|f'\|_\infty$ 完备，而 $\|f\|_\infty$ 不完备；恒等映射从前者到后者连续双射，逆不连续，因为 $f_n(t)=\sin(nt)/n$ 一致趋零，导数却不一致趋零。
\end{example}

\begin{exercise}
设 $T:X\to Y$ 为 Banach 空间之间的有界单射。证明：$R(T)$ 闭当且仅当存在 $c>0$，使 $\|Tx\|\ge c\|x\|$ 对所有 $x$ 成立。
\end{exercise}
\begin{solution}
若值域闭，$R(T)$ 按 $Y$ 的范数完备，有界逆定理用于 $T:X\to R(T)$，得所需下界。反之，若 $Tx_n\to y$，下界使 $x_n$ 为 Cauchy 序列，源空间完备性给出 $x_n\to x$。连续性使 $Tx=y$，故值域闭。单射本身不足，例如 $T:\ell^2\to\ell^2$，$(Tx)_j=x_j/j$；有限支集序列在值域内，但 $(1/j)_j$ 是值域的极限点而不在值域内。
\end{solution}

\section{闭图像定理}\label{sec:yosida-02-06}
\noindent\textit{原书位置：中译本第 79--81 页；英文版第 77--79 页。}

\begin{definition}[闭算子与可闭算子]
线性算子 $T:D(T)\subset X\to Y$ 的图像是 $G(T)=\{(x,Tx):x\in D(T)\}$。若它在 $X\times Y$ 中闭，称 $T$ 闭；若 $\overline{G(T)}$ 仍为某线性算子的图像，称 $T$ 可闭，该算子记为 $\overline T$。对可度量空间，闭性等价于：$x_n\in D(T)$、$x_n\to x$、$Tx_n\to y$ 蕴含 $x\in D(T)$ 且 $Tx=y$。
\end{definition}

\begin{proposition}[可闭判别与逆算子的闭性]\label{prop:y2-closable}
在可度量拓扑向量空间之间，$T$ 可闭当且仅当 $x_n\in D(T)$、$x_n\to0$、$Tx_n\to y$ 蕴含 $y=0$。闭单射的逆 $T^{-1}:R(T)\to X$ 也是闭算子。
\end{proposition}
\begin{proof}
若图像闭包为图像，其中 $(0,y)$ 只能有 $y=0$。反向，$\overline{G(T)}$ 是闭线性子空间，若同时含 $(x,y_1),(x,y_2)$，相减得 $(0,y_1-y_2)$。可度量性允许用图像中的序列逼近它，假设即给出 $y_1=y_2$。所以闭包确为某线性算子的图像，且是包含 $G(T)$ 的最小闭图像。逆算子的图像由同胚 $(x,y)\mapsto(y,x)$ 得到，所以仍闭。
\end{proof}

\begin{theorem}[Banach 闭图像定理]\label{thm:banach-closed-graph}
设 $X,Y$ 为 $F$ 空间，全定义线性算子 $T:X\to Y$ 若闭，则连续。对 Banach 空间，结论等价于存在 $C$ 使 $\|Tx\|\le C\|x\|$。
\end{theorem}
\begin{proof}
积空间 $X\times Y$ 由乘积度量完备，$G(T)$ 是其闭线性子空间，因而也是 $F$ 空间。投影 $\pi_X:G(T)\to X$ 连续线性双射，因为 $D(T)=X$。由开映射定理其逆连续，而 $T=\pi_Y\pi_X^{-1}$ 连续。
\end{proof}

\begin{proposition}[图像范数与算子比较]\label{prop:y2-graph-comparison}
若 $X,Y_1,Y_2$ 为 Banach 空间，$T_1:D(T_1)\subset X\to Y_1$ 闭，则
\[
 \|x\|_{T_1}=\|x\|_X+\|T_1x\|_{Y_1}
\]
使 $D(T_1)$ 成为 Banach 空间。若 $T_2:D(T_2)\subset X\to Y_2$ 可闭且 $D(T_1)\subset D(T_2)$，则存在 $C>0$，使
\[
 \|T_2x\|_{Y_2}\le C(\|x\|_X+\|T_1x\|_{Y_1})\quad(x\in D(T_1)).
\]
\end{proposition}
\begin{proof}
图像范数的 Cauchy 序列同时使 $x_n$ 和 $T_1x_n$ 收敛；闭性保证两个极限仍构成图像中的点，所以范数完备。对 $T_2:D(T_1)\to Y_2$，若 $x_n\to x$ 于图像范数且 $T_2x_n\to y$，则 $x\in D(T_2)$，并有 $x_n-x\to0$、$T_2(x_n-x)\to y-T_2x$。可闭判别给出 $y=T_2x$。故该映射图像闭，由闭图像定理有界。这是原书归于 H\"ormander 的算子比较定理；和范数与平方和范数等价。
\end{proof}

\begin{example}[闭而无界的微分算子]
在 $X=Y=C[0,1]$ 上取 $D(T)=C^1[0,1]$、$Tf=f'$。若 $f_n\to f$、$f_n'\to g$ 都一致，恒等式
\[
 f_n(t)-f_n(0)=\int_0^tf_n'(s)\,ds
\]
传到极限得 $f(t)-f(0)=\int_0^tg(s)\,ds$，故 $f\in C^1$ 且 $f'=g$，所以 $T$ 闭。另一方面 $f_n(t)=t^n$ 满足 $\|f_n\|_\infty=1$、$\|Tf_n\|_\infty=n$，故按 $C$ 范数无界。它的定义域不是整个 $C[0,1]$；换用图像范数以后它是有界的。
\end{example}

\begin{example}[微分算子的可闭性与形式转置]
设 $P=\sum_{|\alpha|\le k}a_\alpha(x)\partial^\alpha$，$a_\alpha\in C^k(\Omega)$。在 $L^2(\Omega)$ 中取
\[
 D(P)=\{u\in C^k(\Omega)\cap L^2(\Omega):Pu\in L^2(\Omega)\}.
\]
若 $u_n\to0$、$Pu_n\to g$ 于 $L^2$，对 $\varphi\in\Test(\Omega)$ 分部积分得
\[
 \int g\varphi=\lim_n\int(Pu_n)\varphi
 =\lim_n\int u_nP^t\varphi=0,\qquad
 P^t\varphi=\sum_{|\alpha|\le k}(-1)^{|\alpha|}\partial^\alpha(a_\alpha\varphi).
\]
紧支集消去边界项，$P^t\varphi\in L^2$ 允许用 Schwarz 不等式取极限。因此 $g=0$，$P$ 可闭。这里的双线性配对不加复共轭，$P^t$ 是形式转置；Hilbert 伴随的系数则须取共轭。
\end{example}

\begin{exercise}
在 $L^2(\R)$ 上定义 $D(T)=\{u\in L^2:xu\in L^2\}$、$(Tu)(x)=xu(x)$。证明 $T$ 闭而无界，并说明其图像范数空间为什么完备。
\end{exercise}
\begin{solution}
若 $u_n\to u$、$xu_n\to v$ 于 $L^2$，可逐次选取同一子列使两者都几乎处处收敛，于是 $v=xu$ 几乎处处。故 $u\in D(T)$ 且 $Tu=v$，图像闭。取 $u_n=\mathbf1_{[n,n+1]}$，有 $\|u_n\|_2=1$、$\|Tu_n\|_2\ge n$，所以无界。图像范数完备由命题 \ref{prop:y2-graph-comparison}，也可对两部分的 Cauchy 序列用刚才的闭性论证。
\end{solution}

\section{闭图像定理的一个应用（H\"ormander 定理）}\label{sec:yosida-02-07}
\noindent\textit{原书位置：中译本第 82--83 页；英文版第 80--81 页。}

本节展示一个有代表性的思路：先由定性光滑性保证导数算子处处有定义，再以闭图像定理得到统一估计，最后代入指数函数，把估计变成多项式零点的限制。

\begin{definition}[本节使用的亚椭圆性]
令 $D_j=(1/i)\partial_{x_j}$，$P(D)$ 为常系数偏微分算子，$P(\zeta)$ 为对应复多项式。本节沿用原书的 $L^2_{\Loc}$ 表述：若在任意开子域上，$u\in L^2_{\Loc}$ 且 $P(D)u$ 在某开集光滑，必有 $u$ 在该开集具有光滑代表，则称 $P(D)$ 亚椭圆。现代通常要求同一结论对所有分布 $u$ 成立；该较强定义当然蕴含本节所用性质。下面的必要条件仅需 $L^2_{\Loc}$ 形式。
\end{definition}

分布等式使用形式转置 $P^t(D)=P(-D)$：$P(D)u=0$ 指 $\int uP(-D)\varphi=0$ 对所有测试函数成立；不加入 Hilbert 配对的复共轭。

\begin{theorem}[H\"ormander 的复特征根必要条件]\label{thm:y2-hormander}
若 $P(D)$ 亚椭圆，则对每个 $A>0$ 存在 $B_A<\infty$，使
\[
 P(\zeta)=0,\quad\zeta=\xi+i\eta\in\mathbb C^n,\quad
 |\eta|\le A\quad\Longrightarrow\quad|\zeta|\le B_A.
\]
等价地，零点集上的 $|\zeta|\to\infty$ 必迫使 $|\operatorname{Im}\zeta|\to\infty$。这采用中译本第 82 页校注中修正的常数次序。
\end{theorem}
\begin{proof}
选非空有界开集 $V\Subset W\subset\R^n$，并使 $W$ 有界。令
\[
 U=\{u\in L^2(W):P(D)u=0\text{ 于分布意义}\}.
\]
对每个测试函数，$u\mapsto\int_WuP(-D)\varphi$ 在 $L^2(W)$ 上连续，故其核闭；$U$ 为这些核的交，所以是 Banach 空间。亚椭圆性保证每个 $u\in U$ 在 $W$ 光滑，从而 $\partial_ju|_V\in L^2(V)$。

算子 $A_j:U\to L^2(V)$、$A_ju=\partial_ju|_V$ 图像闭：若 $u_l\to u$ 于 $L^2(W)$、$A_ju_l\to v$ 于 $L^2(V)$，对 $\varphi\in\Test(V)$ 分部积分取极限得
\[
 \int_Vv\varphi=-\int_Vu\partial_j\varphi,
\]
所以 $v=\partial_ju$。闭图像定理给出常数 $C$，使
\[
 \sum_{j=1}^n\|\partial_ju\|_{L^2(V)}^2
 \le C\|u\|_{L^2(W)}^2\quad(u\in U).
\]
若 $P(\zeta)=0$，取 $u(x)=e^{ix\cdot\zeta}$。因 $W$ 有界，$u\in L^2(W)$ 且 $P(D)u=P(\zeta)u=0$，代入得
\[
 |\zeta|^2\int_Ve^{-2x\cdot\eta}\,dx
 \le C\int_We^{-2x\cdot\eta}\,dx.
\]
选择 $R$ 使 $W\subset B(0,R)$。若 $|\eta|\le A$，左边积分不小于 $|V|e^{-2RA}$，右边积分不大于 $|W|e^{2RA}$，所以
\[
 |\zeta|^2\le C\frac{|W|}{|V|}e^{4RA}.
\]
右边与具体零点无关，取平方根即得 $B_A$。这补全了原书“虚部有界则根有界”的积分比值步骤。
\end{proof}

\begin{example}[Laplace 算子与遗漏方向的导数]
对 Laplace 算子取 $P(\zeta)=-\sum_j\zeta_j^2$。若 $P(\xi+i\eta)=0$，比较实虚部得 $|\xi|^2=|\eta|^2$、$\xi\cdot\eta=0$，所以 $|\zeta|^2=2|\eta|^2$，满足必要条件。但必要条件本身不构成 Weyl 引理的证明；其充分性留待后续理论。

在 $\R^2$ 上，$P(D)=D_1$ 的零点包含 $(0,t)$，$t\in\R$，虚部恒为零而模无界，因而不亚椭圆。也可直接取 $u(x_1,x_2)=|x_2|$：它局部平方可积且 $D_1u=0$，却不光滑。这里的缺陷来自算子没有控制第二个方向。
\end{example}

\begin{exercise}
对常系数算子 $P(D)$，假定每个局部平方可积的齐次解均光滑。证明对任意多重指标 $\alpha$ 及非空有界开集 $V\Subset W$，存在 $C_{\alpha,V,W}$，使
\[
 \|\partial^\alpha u\|_{L^2(V)}\le C_{\alpha,V,W}\|u\|_{L^2(W)}
 \quad(u\in L^2(W),\ P(D)u=0).
\]
这些常数能否由本章论证直接算出？
\end{exercise}
\begin{solution}
齐次解空间 $U$ 在 $L^2(W)$ 中闭。光滑性使 $u\mapsto\partial^\alpha u|_V$ 全定义。若 $u_l\to u$ 于 $U$ 且 $\partial^\alpha u_l\to v$ 于 $L^2(V)$，对测试函数配对得 $\int v\varphi=(-1)^{|\alpha|}\int u\partial^\alpha\varphi$，故 $v=\partial^\alpha u$。闭图像定理给出估计。该论证保证常数存在，并未提供可计算的最佳常数或其随 $\alpha$、区域变化的定量规律。
\end{solution}


\section*{本章参考文献与阅读衔接}
\phantomsection
\addcontentsline{toc}{section}{本章参考文献与阅读衔接}
本章直接核对的资料是两版吉田第二章 \cite{yosida-zh,yosida-en}。原书章末列出 Banach、Bourbaki、Dunford--Schwartz、Hille--Phillips 与 H\"ormander 的著作；它们是原书的进一步阅读线索，本稿未将其列作已直接查阅的文献。中译本第 73 页关于复测度及绝对连续性的校注、第 82 页关于常数次序的校注，已在相应节说明。

本章八道习题均为补充学习题，解答统一输出于书后。第三章转向 Hilbert 空间的投影与表示；分布极限与闭微分算子的例子可与第 \ref{ch:distributions}、\ref{ch:operations}、\ref{ch:pde} 章对读。H\"ormander 特征根条件的充分性和一般 Weyl 引理不在本章证明范围内。
